%%%%%%%%%%%%%%%%%%%%%%%%%%%%%%%%%%%%%%%%%%%%%%%%%%%%%%%%%%%%%%%%%%%%%
%% sec_implementation.tex
%% Section 3: Implementation in RMG
%% RMG Periodic RT-TDDFT -- Paper 2
%%%%%%%%%%%%%%%%%%%%%%%%%%%%%%%%%%%%%%%%%%%%%%%%%%%%%%%%%%%%%%%%%%%%%

The formal theory in
Sec.~\ref{sec:periodic-velocity-gauge}--\ref{sec:optical-response}
defines the periodic velocity-gauge RT-TDDFT formulation and the
corresponding physical response functions.  This section describes the
current numerical realization in \texttt{RMG}.  The implementation
inherits the density-matrix propagation strategy of the previous
finite-system RT-TDDFT module,\cite{jakowski-rmg-tddft-2025} and extends
it to periodic boundary conditions, k-point sampling, collinear spin
polarization, and current-based response analysis.

%--------------------------------------------------------------------
%\subsection{Numerical Implementation and Density-Matrix Propagation}
\subsection{Density-Matrix Propagation}
\label{sec:implementation-density-matrix}
%--------------------------------------------------------------------

A periodic RT-TDDFT calculation begins from a converged ground-state
Kohn--Sham calculation.  For each sampled $\mathbf k$ point and
collinear spin channel $\sigma$, the ground-state cell-periodic Bloch
orbitals are retained as a fixed basis for the real-time propagation.
In this basis, RMG constructs and propagates finite active-space
Hamiltonian and density matrices,
$\mathbf H_{\mathbf k}^{\sigma}(t)$ and
$\mathbf P_{\mathbf k}^{\sigma}(t)$, as introduced formally in
Sec.~\ref{sec:density-matrix-propagation}.  General $\mathbf k$ points
lead to complex Hermitian matrices, while $\Gamma$-point calculations
may use real-valued paths where applicable.

The implementation equations are evaluated internally in atomic units,
with $\hbar=m_e=e=1$ and $q_e=-1$.  The speed of light is not set to
unity by this convention.  The relation between the explicit
electromagnetic quantities in Sec.~\ref{sec:periodic-velocity-gauge}
and the numerical coefficient used by the \texttt{VECTOR\_POT} mode is
therefore treated separately in
Sec.~\ref{sec:implementation-perturbation-current}.  The present
periodic RT-TDDFT path is restricted to norm-conserving
pseudopotentials.

The propagated orbital subspace is selected from the ground-state
Kohn--Sham states.  Its lower boundary is specified by the input
parameter \texttt{tddft\_start\_state}; states below this boundary are
not included in the propagated density matrix.  Let $\{n_a\}$ denote
the ground-state band indices retained in the active subspace and let
$N_{\rm prop}$ be the number of propagated states.  At the beginning of
the real-time calculation,
\begin{equation}
  H_{ab,\mathbf k}^{\sigma}(0)
  =
  \epsilon_{n_a\mathbf k}^{\sigma}\delta_{ab},
  \qquad
  P_{ab,\mathbf k}^{\sigma}(0)
  =
  f_{n_a\mathbf k}^{\sigma}\delta_{ab},
  \label{eq:impl_initial_h_p}
\end{equation}
where $\epsilon_{n_a\mathbf k}^{\sigma}$ and
$f_{n_a\mathbf k}^{\sigma}$ are the ground-state eigenvalue and
occupation of state $n_a$.  Occupied states below the active lower
boundary remain frozen at their ground-state density contribution.
Occupied and unoccupied states inside the active subspace are included
in $\mathbf P_{\mathbf k}^{\sigma}(t)$ and can participate in the
real-time response.

For a time step $\Delta t$, RMG uses the first term of the Magnus
expansion with an endpoint-averaged Hamiltonian.  In atomic units, the
mathematically equivalent generator is
\begin{equation}
  \boldsymbol{\Omega}_{\mathbf k}^{\sigma}(t+\Delta t,t)
  \simeq
  -\frac{\imath\Delta t}{2}
  \left[
    \mathbf H_{\mathbf k}^{\sigma}(t)
    +
    \mathbf H_{\mathbf k}^{\sigma}(t+\Delta t)
  \right].
  \label{eq:impl_endpoint_magnus}
\end{equation}
The code stores the Hermitian Hamiltonian integral separately from the
factor of $-\imath$ used in the commutator propagation, but the
resulting update is equivalent to
\begin{equation}
  \mathbf P_{\mathbf k}^{\sigma}(t+\Delta t)
  =
  e^{\boldsymbol{\Omega}_{\mathbf k}^{\sigma}}
  \mathbf P_{\mathbf k}^{\sigma}(t)
  e^{-\boldsymbol{\Omega}_{\mathbf k}^{\sigma}} .
  \label{eq:impl_density_update}
\end{equation}
This unitary transformation is applied through the BCH
nested-commutator expansion.  The complex periodic implementation uses
this BCH propagation route.

The endpoint Hamiltonian
$\mathbf H_{\mathbf k}^{\sigma}(t+\Delta t)$ is obtained
self-consistently.  RMG first predicts the endpoint Hamiltonian from
the preceding Hamiltonian history.  With this trial endpoint,
Eq.~\eqref{eq:impl_endpoint_magnus} is used to propagate the density
matrix, the real-space spin density is reconstructed, and the Hartree
and spin-dependent exchange--correlation potentials are updated.  The
matrix elements of the resulting change in the local effective
potential are then added to the endpoint Hamiltonian, and the cycle is
repeated until the endpoint Hamiltonian is self consistent:
\[
  \mathbf H(t+\Delta t)
  \rightarrow
  \boldsymbol{\Omega}
  \rightarrow
  \mathbf P(t+\Delta t)
  \rightarrow
  \rho^\sigma(\mathbf r,t+\Delta t)
  \rightarrow
  V_{\rm H}[\rho],
  V_{\rm xc}^{\sigma}[\rho^\uparrow,\rho^\downarrow]
  \rightarrow
  \mathbf H(t+\Delta t).
\]
The two collinear spin channels are propagated as separate density
matrices, but the potentials are coupled through the total density
entering the Hartree term and through the spin dependence of the
exchange--correlation functional.

The real-space density reconstruction follows the active-space form of
Sec.~\ref{sec:density-matrix-propagation}.  RMG forms the active
density-matrix change
\begin{equation}
  \Delta P_{ab,\mathbf k}^{\sigma}(t)
  =
  P_{ab,\mathbf k}^{\sigma}(t)
  -
  f_{n_a\mathbf k}^{\sigma}\delta_{ab},
  \label{eq:impl_delta_p_active}
\end{equation}
and reconstructs
\begin{equation}
  \rho^\sigma(\mathbf r,t)
  =
  \rho_{\rm gs}^{\sigma}(\mathbf r)
  +
  \sum_{\mathbf k}
  w_{\mathbf k}
  \sum_{ab=1}^{N_{\rm prop}}
  u_{n_a\mathbf k}^{\sigma}(\mathbf r)
  \Delta P_{ab,\mathbf k}^{\sigma}(t)
  u_{n_b\mathbf k}^{\sigma *}(\mathbf r).
  \label{eq:impl_density_reconstruction}
\end{equation}
The ordering of the orbital factors in
Eq.~\eqref{eq:impl_density_reconstruction} matches the density-matrix
convention of Sec.~\ref{sec:density-matrix-propagation}: the
ground-state density supplies the frozen occupied contribution and the
reference occupation of the active subspace, while
$\Delta\mathbf P_{\mathbf k}^{\sigma}(t)$ supplies the time-dependent
active-space density change.  This construction makes convergence with
respect to the number of included virtual states an explicit numerical
parameter of the RT-TDDFT calculation.

Exact unitary propagation preserves the trace, Hermiticity, and
eigenvalues of each density matrix.  In practical calculations these
properties are maintained up to errors from the Magnus truncation, the
finite BCH expansion, and the convergence tolerance of the endpoint
self-consistency cycle.

%--------------------------------------------------------------------
\subsection{External Perturbation and Current Evaluation}
\label{sec:implementation-perturbation-current}
%--------------------------------------------------------------------

In the formal velocity-gauge theory, a prescribed homogeneous vector
potential $\mathbf A(t)$ remains in the Hamiltonian for all times
required by the external field.  The present RMG \texttt{VECTOR\_POT}
mode implements a more specific discrete momentum-coupled excitation
used to initiate the real-time response.  We denote by
$\kappa_\alpha$ the Cartesian numerical coefficient obtained from the
\texttt{electric\_field\_tddft} input after the code's coordinate
conversion.  This symbol denotes the coefficient actually used by the
implementation; it is not identified here uniquely with a physical
electric-field amplitude, vector-potential amplitude, $A_\alpha/c$, or
$eA_\alpha/c$.

Before applying the perturbation, RMG constructs the kinetic/Bloch
momentum-like matrix
\begin{equation}
  M_{ab,\mathbf k,\alpha}^{\rm kin,\sigma}
  =
  \imath
  \left\langle
    \left(\nabla_\alpha+\imath k_\alpha\right)
    u_{n_a\mathbf k}^{\sigma}
    \middle|
    u_{n_b\mathbf k}^{\sigma}
  \right\rangle_{\rm cell}.
  \label{eq:impl_mkin}
\end{equation}
Using integration by parts under periodic boundary conditions, this is
equivalent to the cell-periodic matrix element of
$-\imath\nabla_\alpha+k_\alpha$ in atomic units.  The explicit
$\mathbf k$ term is the Bloch contribution to the momentum operator.

The implemented initial perturbation is then
\begin{equation}
  \Delta\mathbf H_{\mathbf k}^{\rm code,\sigma}
  =
  \sum_\alpha
  \kappa_\alpha
  \mathbf M_{\mathbf k,\alpha}^{\rm kin,\sigma}.
  \label{eq:impl_vector_pot_kick}
\end{equation}
The verified construction order is
\[
  \texttt{VecPHmatrix()}
  \rightarrow
  \texttt{VECTOR\_POT\ perturbation}
  \rightarrow
  \texttt{CurrentNlpp()} .
\]
Consequently, the perturbing operator in
Eq.~\eqref{eq:impl_vector_pot_kick} contains the kinetic/Bloch
contribution only.  The nonlocal-pseudopotential commutator correction
does not enter the initial perturbing kick in the present
implementation.

The perturbing contribution enters the first accepted midpoint
propagator through the initial Hamiltonian endpoint and is absent from
the converged self-consistent Hamiltonian without the external
perturbing term at subsequent times.  Since the first Magnus generator
uses an endpoint-averaged Hamiltonian, the effective integrated
first-order numerical action of the perturbation is
\begin{equation}
  \int dt\,
  \Delta\mathbf H_{\mathbf k}^{\rm code,\sigma}(t)
  \approx
  \frac{\Delta t}{2}
  \sum_\alpha
  \kappa_\alpha
  \mathbf M_{\mathbf k,\alpha}^{\rm kin,\sigma}.
  \label{eq:impl_first_step_action}
\end{equation}
Equation~\eqref{eq:impl_first_step_action} describes the current
discrete numerical realization.  It should not be read as a statement
that the code applies a mathematically exact Dirac-delta vector
potential or the formally persistent step-like vector potential of
Sec.~\ref{sec:optical-response}.

After the perturbation is applied, RMG adds the nonlocal
pseudopotential contribution to the matrices used for current
evaluation.  The routine \texttt{CurrentNlpp()} calls the same
directional nonlocal-operator path used for the projector commutator.
Expressed in the sign convention of Sec.~\ref{sec:current-nlpp}, the
atomic-unit nonlocal velocity contribution is
$\imath[\hat V_{\rm NL},\hat r_\alpha]$, or equivalently
$(1/\imath)[\hat r_\alpha,\hat V_{\rm NL}]$.  The corresponding raw
matrix contribution is written as
\begin{equation}
  M_{ab,\mathbf k,\alpha}^{\rm NL,\sigma}
  =
  \frac{1}{\imath}
  \left\langle
    u_{n_a\mathbf k}^{\sigma}
    \left|
    \left[
      \hat r_\alpha,\hat V_{{\rm NL},\mathbf k}^{\sigma}
    \right]
    \right|
    u_{n_b\mathbf k}^{\sigma}
  \right\rangle_{\rm cell}.
  \label{eq:impl_mnl}
\end{equation}
Equivalently, this term may be written as
$\imath\langle u_{n_a\mathbf k}^{\sigma}|
[\hat V_{{\rm NL},\mathbf k}^{\sigma},\hat r_\alpha]
|u_{n_b\mathbf k}^{\sigma}\rangle_{\rm cell}$.
This is the same sign convention used in the formal current matrix of
Eq.~\eqref{eq:bloch_current_matrix}, excluding the charge factor
$q_e$ and the volume normalization that are not included in the raw
RMG matrix.

Thus the raw current-evaluation matrix used by RMG is
\begin{equation}
  \mathbf M_{\mathbf k,\alpha}^{\rm raw,\sigma}
  =
  \mathbf M_{\mathbf k,\alpha}^{\rm kin,\sigma}
  +
  \mathbf M_{\mathbf k,\alpha}^{\rm NL,\sigma}.
  \label{eq:impl_mraw}
\end{equation}
For calculations using the simple DFT+U implementation, the same
directional operator-construction path also includes the corresponding
Hubbard-projector contribution,
\begin{equation}
  \mathbf M_{\mathbf k,\alpha}^{\rm raw,\sigma}
  =
  \mathbf M_{\mathbf k,\alpha}^{\rm kin,\sigma}
  +
  \mathbf M_{\mathbf k,\alpha}^{\rm NL,\sigma}
  +
  \mathbf M_{\mathbf k,\alpha}^{U,\sigma}.
  \label{eq:impl_mraw_ldau}
\end{equation}
This statement identifies the implemented code path; it does not by
itself establish that the DFT+U contribution has been independently
proven to be the full gauge-consistent commutator for a Hubbard
functional.

The distinction between perturbation and measurement is therefore
important.  The current RMG \texttt{VECTOR\_POT} perturbation uses the
kinetic/Bloch momentum-like matrix in
Eq.~\eqref{eq:impl_vector_pot_kick}, whereas the matrix subsequently
used to evaluate the current-like response contains the nonlocal
pseudopotential correction and, where active, the simple DFT+U
projector contribution.  The implementation also does not explicitly
construct the $A$-dependent diamagnetic current operator appearing in
the formal charge-current operator of Sec.~\ref{sec:current-nlpp}.

%--------------------------------------------------------------------
\subsection{Frequency Response of the Impulsive RMG Velocity-Gauge Perturbation}
\label{sec:rmg_vg_frequency_response}

The temporal form of the external perturbation determines the frequency
dependence that must be removed from the calculated current in order to
recover a material response function.  This distinction is particularly
important for the present RMG implementation.

For reference, consider first the conventional velocity-gauge
representation of an impulsive electric field,

\begin{equation}
E_{\beta}(t)=E_{0,\beta}\delta(t).
\label{eq:standard_e_delta}
\end{equation}

As shown in Sec.~\ref{sec:optical_response}, the gauge-equivalent vector
potential is a persistent step,

\begin{equation}
A_{\beta}(t)
=
-cE_{0,\beta}\Theta(t).
\label{eq:standard_A_step}
\end{equation}

Because $E_{\beta}(\omega)=E_{0,\beta}$ is independent of frequency,
the corresponding optical conductivity follows directly from the
Fourier-transformed current,

\begin{equation}
\sigma_{\alpha\beta}(\omega)
=
\frac{\Delta J_{\alpha}^{(\beta)}(\omega)}
     {E_{0,\beta}}.
\label{eq:standard_vg_conductivity}
\end{equation}

The impulsive perturbation used in the present RMG implementation is
different.  Its idealized continuous-time form may be represented as a
vector-potential impulse,

\begin{equation}
A_{\beta}^{\mathrm{imp}}(t)
=
\mathcal{A}_{0,\beta}\delta(t),
\label{eq:A_delta_impulse}
\end{equation}

where $\mathcal{A}_{0,\beta}$ denotes the integrated strength of the
vector-potential-like impulse.  Since

\begin{equation}
E_{\beta}(t)
=
-\frac{1}{c}
\frac{dA_{\beta}(t)}{dt},
\label{eq:E_from_A_again}
\end{equation}

the corresponding electric field is

\begin{equation}
E_{\beta}^{\mathrm{imp}}(t)
=
-\frac{\mathcal{A}_{0,\beta}}{c}
\delta'(t).
\label{eq:E_delta_prime}
\end{equation}

Using the Fourier-transform identity

\begin{equation}
\mathcal{F}
\left[
\frac{df(t)}{dt}
\right]
=
i\omega f(\omega),
\label{eq:fft_derivative_repeated}
\end{equation}

one obtains

\begin{equation}
E_{\beta}^{\mathrm{imp}}(\omega)
=
-\frac{i\omega}{c}
\mathcal{A}_{0,\beta}.
\label{eq:E_impulse_frequency}
\end{equation}

Thus, unlike the conventional velocity-gauge step perturbation, whose
electric-field spectrum is frequency independent, the electric field
associated with the impulsive vector potential carries an additional
factor of $i\omega$.

In linear response,

\begin{equation}
\Delta J_{\alpha}(\omega)
=
\sum_{\beta}
\sigma_{\alpha\beta}(\omega)
E_{\beta}(\omega),
\label{eq:J_sigma_E_rmg}
\end{equation}

so that for the impulsive RMG perturbation,

\begin{equation}
\sigma_{\alpha\beta}^{\mathrm{RMG}}(\omega)
=
-\frac{c}
       {i\omega\mathcal{A}_{0,\beta}}
\Delta J_{\alpha}^{(\beta)}(\omega).
\label{eq:rmg_sigma_from_current}
\end{equation}

Apart from the overall perturbation amplitude and unit conventions,
Eq.~\eqref{eq:rmg_sigma_from_current} shows that reconstruction of the
conductivity from the present RMG current response requires one factor
of $1/(i\omega)$.  This factor originates from the temporal form of the
applied perturbation; it is not the current-to-polarization conversion.

A second and physically distinct factor of $1/(i\omega)$ appears when
the conductivity is converted to polarization or susceptibility.  From

\begin{equation}
\Delta J_{\alpha}(\omega)
=
i\omega\Delta P_{\alpha}(\omega)
\end{equation}

and

\begin{equation}
\Delta P_{\alpha}(\omega)
=
\sum_{\beta}
\chi_{\alpha\beta}(\omega)
E_{\beta}(\omega),
\end{equation}

one obtains

\begin{equation}
\sigma_{\alpha\beta}(\omega)
=
i\omega
\chi_{\alpha\beta}(\omega),
\qquad
\chi_{\alpha\beta}(\omega)
=
\frac{\sigma_{\alpha\beta}(\omega)}
     {i\omega}.
\label{eq:sigma_chi_relation_rmg}
\end{equation}

Combining Eqs.~\eqref{eq:rmg_sigma_from_current} and
\eqref{eq:sigma_chi_relation_rmg} therefore gives

\begin{equation}
\chi_{\alpha\beta}^{\mathrm{RMG}}(\omega)
=
\frac{c}
     {\mathcal{A}_{0,\beta}}
\frac{\Delta J_{\alpha}^{(\beta)}(\omega)}
     {\omega^{2}},
\label{eq:rmg_chi_current_omega2}
\end{equation}

up to the overall sign and normalization conventions associated with
the definition of the numerical perturbation and current.

The two factors of $1/\omega$ in
Eq.~\eqref{eq:rmg_chi_current_omega2} therefore have different physical
origins.  The first compensates the additional $i\omega$ frequency
dependence of the electric field generated by the impulsive
vector-potential perturbation, while the second converts the
conductivity to a polarization or susceptibility response.  The same
analysis also predicts a phase change associated with the factors of
$i$, so comparisons of complex response functions require consistent
treatment of both their real and imaginary components.

The actual RMG perturbation is implemented over a finite numerical
propagation interval rather than as a mathematical Dirac delta
function.  The correspondence between the idealized impulse of
Eq.~\eqref{eq:A_delta_impulse} and the finite-$\Delta t$ RMG
implementation is discussed in Appendix~\ref{app:finite_dt_impulse}.

%--------------------------------------------------------------------
\subsection{Periodic, Spin-Resolved, and Parallel Implementation}
\label{sec:implementation-periodic-spin-parallel}
%--------------------------------------------------------------------

The periodic implementation propagates a separate active-space
Hamiltonian and density matrix for each local $\mathbf k$ point and
collinear spin channel.  Thus
$\mathbf H_{\mathbf k}^{\sigma}(t)$,
$\mathbf P_{\mathbf k}^{\sigma}(t)$, and the corresponding
momentum/current-like matrices are stored and updated independently for
each $\mathbf k$ handled by a given MPI task group.  General
$\mathbf k$ points require complex Bloch orbitals and complex
Hermitian matrices, while $\Gamma$-point calculations may use
real-valued orbital and matrix paths where applicable.  The Bloch
$\mathbf k$ contribution to the kinetic momentum matrix is included in
the operator construction described in
Sec.~\ref{sec:implementation-perturbation-current}.

The density and observable reductions use the Brillouin-zone weights
associated with the sampled $\mathbf k$ points.  During the
self-consistent endpoint update, RMG reconstructs a density
contribution from each local $\mathbf k$ point, multiplies it by the
corresponding weight, and sums over the distributed $\mathbf k$ set.
The same weighted reduction is used for current-like vector
observables.  After the k-point and matrix-distribution reductions,
the vector current is passed through the RMG symmetry machinery for
vector symmetrization.

For collinear spin-polarized calculations, RMG propagates the two spin
channels as separate matrix problems.  Each spin channel has its own
$\mathbf P_{\mathbf k}^{\uparrow}$ or
$\mathbf P_{\mathbf k}^{\downarrow}$ and its own spin-dependent
Hamiltonian.  The implementation does not introduce off-diagonal
spinor blocks, spin flips, transverse spin dynamics, or spin-orbit
induced spin mixing in the collinear RT-TDDFT path used here.  The
channels are nevertheless coupled self-consistently: after the
k-weighted density reconstruction, the spin density available on a
given spin partition is combined with the opposite-spin density before
the Hartree and exchange--correlation potentials are evaluated.  The
Hartree update depends on the total density, while the
exchange--correlation update uses the spin-dependent functional form
with both spin densities.

Current-like observables are accumulated separately for each spin
channel in the implementation.  In the notation of
Secs.~\ref{sec:collinear-spin-extension} and~\ref{sec:optical-response},
the physical total charge-current response is obtained from the
spin-channel contributions,
\begin{equation}
  \Delta J_\alpha(t)
  =
  \Delta J_\alpha^\uparrow(t)
  +
  \Delta J_\alpha^\downarrow(t).
  \label{eq:impl_spin_sum_current}
\end{equation}
The separate spin-channel traces can also be retained for
spin-resolved analysis within the fixed collinear basis.  This
organization is used later to analyze the spin-up and spin-down
optical channels of the NV$^{-}$ center.

The parallel decomposition combines the real-space domain
decomposition inherited from ground-state RMG with additional
parallelism over k points, spin channels, and dense active-space
matrices.  Real-space operations such as density reconstruction,
application of local potentials, and grid integrations are distributed
over the real-space grid.  The active-space Hamiltonian and density
matrices are distributed using the matrix backend selected for the
run, including ScaLAPACK and tiled matrix-multiplication paths.  The
refactored \texttt{rmg::tddft<OrbitalType,MatrixType>} class provides
CPU and GPU execution paths through templated orbital and matrix
types.

This organization exposes independent work over $\mathbf k$ points and
spin channels while retaining the locality of the real-space RMG
operations and the efficiency of distributed dense matrix algebra.  It
is this combination of real-space decomposition, k/spin decomposition,
matrix parallelism, and accelerator support that makes the periodic
density-matrix RT-TDDFT implementation applicable to large supercells,
including the 4095-atom NV$^{-}$-center calculations discussed below.

%--------------------------------------------------------------------
\subsection{Optical Spectra}% and Implementation-Specific Processing}
\label{sec:implementation-optical-processing}
%--------------------------------------------------------------------

The formal response relations in Sec.~\ref{sec:optical-response}
define the optical conductivity and dielectric tensor in terms of the
physical macroscopic charge-current density.  RMG writes
spin-resolved raw current-like traces, for which the available
current-postprocessing workflow produces optical spectra in arbitrary
units.

Using the raw current-evaluation matrix defined in
Sec.~\ref{sec:implementation-perturbation-current}, the quantity
written by RMG for spin channel $\sigma$ can be represented
schematically as
\begin{equation}
  R_\alpha^\sigma(t_n)
  =
  {\cal S}_\alpha
  \left\{
  \sum_{\mathbf k}
  w_{\mathbf k}\,
  {\rm Re}\,
  {\rm Tr}
  \left[
    \mathbf P_{\mathbf k}^{\sigma}(t_n)
    \mathbf M_{\mathbf k,\alpha}^{\rm raw,\sigma}
  \right]
  \right\},
  \label{eq:impl_raw_current_trace}
\end{equation}
where ${\cal S}$ denotes the vector symmetrization applied after the
parallel reductions.  The raw trace includes k-point weights and the
reduction over distributed k-point and matrix partitions.  It is
written separately for each spin channel and does not explicitly
include multiplication by $q_e$, division by $V_{\rm cell}$, or the
sum over spin channels.

Consequently, $R_\alpha^\sigma(t)$ should not be identified directly
with the physical macroscopic charge-current density
$J_\alpha^\sigma(t)$ of Sec.~\ref{sec:current-nlpp}.  If the raw
matrix is interpreted as a zero-field velocity-like matrix in atomic
units, conversion to the formal current density would require, at
minimum, the electron-charge factor, the cell-volume normalization,
the spin summation for the total current, and the corresponding unit
conversion:
\begin{equation}
  J_\alpha(t)
  \sim
  \frac{q_e}{V_{\rm cell}}
  \sum_{\sigma=\uparrow,\downarrow}
  R_\alpha^\sigma(t),
  \label{eq:impl_raw_to_current_schematic}
\end{equation}
with $q_e=-1$ in the internal atomic-unit convention.  This is a
schematic relation identifying missing normalization factors, not an
operation currently applied by the RMG current-output routine.  The
code also prints a ground-state current as a diagnostic header entry,
but this value is not directly subtracted from the written time trace.
In a fresh calculation, the first numerical current row is written
after the first accepted propagation step, using the step-counter time
label.

For current-based spectral analysis, the available RMG workflow first
removes a numerical baseline from the sampled trace.  For current
files, it then applies a cumulative summation of the current-like
signal, which acts as an implementation-specific discrete
current-to-polarization integration.  The resulting time series may be
detrended and damped or windowed, Fourier transformed, normalized by
the raw perturbation coefficient, and smoothed to produce the plotted
spectrum.  This processing is useful for comparing spectral shapes and
peak positions, but the available workflow does not preserve the full
signed complex phase of the response.  It therefore does not by itself
produce a quantitatively normalized complex conductivity or dielectric
function.

The spectra presented below are reported in arbitrary units and are
interpreted primarily through relative peak positions, spectral shape,
polarization dependence, spin dependence, and active-space convergence
rather than absolute optical intensities.  A fully normalized
conductivity or dielectric response can be obtained from the formal
relations of Sec.~\ref{sec:optical-response}, but doing so requires
retaining the complex Fourier response and applying the physical
charge, volume, spin-summation, field-normalization, and unit-conversion
factors explicitly.
